% Self-contained local starter inspired by the NTNU Laboratory Report in Physics
% template by Vebjorn S. Forde (CC BY 4.0):
% https://www.overleaf.com/latex/templates/ntnu-laboratory-report-in-physics/zrkxstqtpdpv
% This is a simplified adaptation, not the original multi-file project.
% Compile with pdfLaTeX in TeXstudio. Replace the example text and metadata.

\documentclass[5p,a4paper]{elsarticle}

% Common TeX Live packages
\usepackage[utf8]{inputenc}
\usepackage[T1]{fontenc}
\usepackage[english]{babel}
\usepackage{lmodern}
\usepackage{geometry}
\geometry{margin=25mm}
\usepackage{microtype}
\usepackage{graphicx}
\usepackage{amsmath,amssymb}
\usepackage{siunitx}
\usepackage{booktabs}
\usepackage{float}
\usepackage{caption}
\usepackage{enumitem}
\usepackage[hidelinks]{hyperref}
\usepackage[numbers]{natbib}

% ---------- Report details: edit these ----------
\newcommand{\reporttitle}{Title of the Laboratory Report}
\newcommand{\course}{Course name and code}
\newcommand{\studentname}{Your name}
\newcommand{\studentid}{Student number (optional)}
\newcommand{\labdate}{Date of experiment}
\newcommand{\submissiondate}{Date of submission}
\newcommand{\institution}{Department / University}

\title{\reporttitle}
\author{\studentname\\\small \studentid\\\small \course\\\small \institution}
\date{\small Experiment: \labdate \quad\textbar\quad Submitted: \submissiondate}

\begin{document}
	\maketitle
	
	\begin{abstract}
		Briefly state the purpose of the experiment, the method, the most important
		quantitative result, and the main conclusion. Keep this self-contained and
		usually short. Example: The acceleration due to gravity was estimated from
		measurements of a falling object to be \SI{9.8}{\metre\per\second\squared},
		within the measurement uncertainty.
	\end{abstract}
	
	\noindent\textbf{Keywords:} experiment; measurement; uncertainty
	
	\section{Introduction}
	Explain the scientific question and why it matters. State the objective of the
	experiment and give enough context for the reader to understand the report.
	Cite sources where appropriate, for example \citep{example-reference}.
	
	\section{Theory}
	Present the principles and equations needed to interpret the measurements.
	Define every symbol and include assumptions. For example, a linear model can be
	written as
	\begin{equation}
		y = ax+b,
		\label{eq:linear-model}
	\end{equation}
	where $a$ is the slope and $b$ is the intercept.
	
	\section{Method and Equipment}
	Describe the apparatus and procedure with enough detail that another person
	could repeat the experiment. Include a labelled figure if it helps. Put image
	files in the same folder as this \texttt{.tex} file and replace the filename
	below when you add one.
	
	% Uncomment and edit this figure block when you have an image:
	% \begin{figure}[H]
		%     \centering
		%     \includegraphics[width=0.7\linewidth]{apparatus.png}
		%     \caption{Labelled experimental setup.}
		%     \label{fig:setup}
		% \end{figure}
	
	\section{Results}
	Report observations and calculated results clearly. Keep interpretation for
	the discussion. Include units and uncertainty, and refer to each table or
	figure in the text. Table~\ref{tab:example} shows a simple example.
	
	\begin{table}[H]
		\centering
		\caption{Example measurements. Replace these values with your data.}
		\label{tab:example}
		\begin{tabular}{ccc}
			\toprule
			Trial & $x$ (unit) & $y$ (unit) \\
			\midrule
			1 & 1.0 & 2.1 \\
			2 & 2.0 & 4.0 \\
			3 & 3.0 & 6.2 \\
			\bottomrule
		\end{tabular}
	\end{table}
	
	\section{Discussion}
	Interpret the results in relation to the theory and objective. Compare with
	expected or reference values when relevant. Discuss uncertainty, likely error
	sources, limitations, and possible improvements. Explain unexpected results
	rather than simply listing them.
	
	\section{Conclusion}
	Summarise the main finding and whether the objective was achieved. Include the
	most important numerical result with its uncertainty when applicable. Do not
	introduce new results here.
	
	\section*{Acknowledgements}
	Optional. Acknowledge help or resources that contributed to the work.
	
	\begin{thebibliography}{9}
		\bibitem{example-reference}
		Author, ``Title of source,'' Publisher or journal, year. Replace this example
		with the sources you actually used.
	\end{thebibliography}
	
	\appendix
	\section{Supporting calculations}
	Optional. Put extended derivations, additional data, or other supporting
	material here. Remove this appendix if it is not needed.
	
\end{document}